\subsection{Fourier Approximation}

\begin{definition} \label{def:periodic-convolution-and-fejer-kernel} \hlt{Periodic Convolution and Fejer Kernels}\\
For $f,g\in C(\R/\Z;\C)$, define their \hlt{periodic convolution} as $(f*g)(x)\defeq\int_0^1 f(y)g(x-y)\,dy$.\\
For $N\geq1$, the \hlt{Fejer kernel} is $F_N(x)\defeq \sum\limits_{n=-N}^{N}\left(1-\frac{\abs{n}}{N}\right)e_n(x) =\frac{1}{N}\abs{\sum\limits_{n=0}^{N-1}e_n(x)}^2$.
\end{definition}

\begin{figure}[H]
  \centering
  \begin{tikzpicture}[
  >=Stealth,
  every node/.style={font=\scriptsize},
  x=5.8cm,
  y=0.48cm,
  axis/.style={->, black!55, line width=0.45pt},
  guide/.style={black!35, dashed, line width=0.45pt},
  broad/.style={black!55, dashed, line width=0.75pt},
  sharp/.style={black!85, line width=0.9pt},
  window/.style={fill=gray!12, draw=none}
]
  \draw[window] (-0.10,0) rectangle (0.10,5.65);

  \draw[axis] (-0.52,0) -- (0.55,0) node[right] {$x$};
  \draw[axis] (0,-0.18) -- (0,5.86);

  \draw[guide] (-0.5,0) -- (-0.5,5.65);
  \draw[guide] (0.5,0) -- (0.5,5.65);
  \node[below] at (-0.5,0) {$-\frac12$};
  \node[below] at (0,0) {$0$};
  \node[below] at (0.5,0) {$\frac12$};

  \draw[broad, domain=-0.5:0.5, samples=180, smooth]
    plot (\x,{2.15*exp(-42*\x*\x)});
  \draw[sharp, domain=-0.5:0.5, samples=220, smooth]
    plot (\x,{5.15*exp(-360*\x*\x)});

  \draw[guide] (-0.10,0) -- (-0.10,5.4);
  \draw[guide] (0.10,0) -- (0.10,5.4);
  \draw[<->, black!55] (-0.10,-0.46) -- node[below] {$2\delta$} (0.10,-0.46);

  \node[black!55] at (0.26,1.38) {$F_M$};
  \node[black!85] at (0.12,4.55) {$F_N$};
\end{tikzpicture}

  \caption{The Fejer kernels are nonnegative trigonometric polynomial approximate identities on $\R/\Z$.}
  \label{fig:fejer-kernel-periodic-approx-identity}
\end{figure}

\begin{theorem} \label{thm:fejer-periodic-approximation-identity} \hlt{Fejer Approximation on $\R/\Z$}\\
For $N\geq1$, $F_N$ is a nonnegative trigonometric polynomial with $\int_0^1F_N=1$. For each $0<\delta<1/2$, $F_N\to0$ uniformly on $[\delta,1-\delta]$. If $f\in C(\R/\Z;\C)$, then $f*F_N$ is a trigonometric polynomial and $f*F_N\rightrightarrows f$.
\begin{proof}
From Definition \ref{def:periodic-convolution-and-fejer-kernel}, $F_N$ is a finite linear combination of $e_n$ and $F_N=\frac1N\abs{\sum_{n=0}^{N-1}e_n}^2\geq0$.\\
Since $\int_0^1e_n=0$ for $n\neq0$, the $e_0$ coefficient gives $\int_0^1F_N=1$.\\
For $x\notin\Z$, $F_N(x)=\frac{\sin^2(\pi Nx)}{N\sin^2(\pi x)}$. If $m_\delta\defeq\min_{\delta\leq x\leq1-\delta}\sin^2(\pi x)>0$, then $0\leq F_N(x)\leq (Nm_\delta)^{-1}$ on $[\delta,1-\delta]$.\\
Also $f*e_n=\widehat f_{\mathbb{T}}(n)e_n$, so $f*F_N$ is a trigonometric polynomial. For uniform convergence, write $(f*F_N)(x)=\int_0^1f(x-t)F_N(t)\,dt$. Let $M=\norm{f}_\infty$ and choose $\delta$ so $\abs{f(x-t)-f(x)}<\varepsilon$ whenever $t\in[0,\delta]\cup[1-\delta,1]$.\\
By $F_N\geq0$ and $\int_0^1F_N=1$, uniformly in $x$, $\abs{(f*F_N)(x)-f(x)} \leq \varepsilon+2M\sup_{\delta\leq t\leq1-\delta}F_N(t)\to \varepsilon$.\\
\hlt{Structural Commentary}: Fejer kernels concentrate mass near the identity element of $\R/\Z$.\\
\hlt{Key Insight}: Positivity plus concentration turns convolution into uniform approximation.
\end{proof}
\end{theorem}

\begin{perspective} \label{per:fourier-plancherel-roadmap-tao} \hlt{Fourier and Plancherel Roadmap}\\
Path: orthonormal characters, Fejer density, best $L^2$ projection, then Parseval/Plancherel.\\
\hlt{Key Insight}: Fourier theory is finite-dimensional orthogonal projection plus an approximation theorem.
\end{perspective}

\begin{perspective} \label{per:fourier-coefficients-as-projections} \hlt{Fourier Coefficients as Projections}\\
Fourier coefficients are projections onto orthonormal directions, modeled finitely by Theorems \ref{thm:orthogonal-direct-sum-decomposition}, \ref{thm:orthonormal-basis-coordinates}, and \ref{thm:bessel-inequality-finite-orthonormal-list}; infinite directions make convergence separate.
\end{perspective}

\begin{theorem} \label{thm:orthonormal-fourier-best-approximation-bessel} \hlt{Fourier Coefficients Give Best Mean-Square Approximation}\\
Let $\langle u,v\rangle\defeq\int_a^b u\overline v$ and $\norm{u}_2^2\defeq\langle u,u\rangle$.\\
If $(\phi_n)$ is orthonormal, $c_m=\langle f,\phi_m\rangle$, $s_n=\sum\limits_{m=1}^{n}c_m\phi_m$, and $t_n=\sum\limits_{m=1}^{n}\gamma_m\phi_m$, then $\norm{f-s_n}_2\leq\norm{f-t_n}_2$, with equality iff $\gamma_m=c_m$ for $m=1,\ldots,n$. In particular, $\sum\limits_{m=1}^{\infty}\abs{c_m}^2\leq\norm{f}_2^2$ for $c_m\to0$.
\begin{proof}
For sums over $m=1,\ldots,n$, orthonormality gives\\
$\norm{f-t_n}_2^2 =\norm{f}_2^2-2\operatorname{Re}\sum c_m\overline{\gamma_m}+\sum\abs{\gamma_m}^2 =\norm{f}_2^2-\sum\abs{c_m}^2+\sum\abs{\gamma_m-c_m}^2$.\\
Thus the minimum occurs exactly when $\gamma_m=c_m$.\\
Taking $t_n=s_n$ gives $\sum_{m=1}^{n}\abs{c_m}^2=\norm{s_n}_2^2\leq\norm{f}_2^2$; let $n\to\infty$ to get Bessel and $c_m\to0$.\\
\hlt{Structural Commentary}: Orthogonality makes the error split into independent coefficient errors.\\
\hlt{Key Insight}: Fourier coefficients are the coordinates of the orthogonal projection.
\end{proof}
\end{theorem}

\begin{theorem} \label{thm:trigonometric-fourier-dirichlet-localization} \hlt{Dirichlet Kernel and Local Fourier Convergence}\\
Let $f$ be Riemann integrable and $2\pi$-periodic, with $\widehat{f}(m)\defeq\frac{1}{2\pi}\int_{-\pi}^{\pi}f(t)e^{-imt}\,dt$.\\
Let $s_N(f;x)\defeq\sum\limits_{m=-N}^{N}\widehat{f}(m)e^{imx}$ and $D_N(t)\defeq\sum\limits_{m=-N}^{N}e^{imt}=\frac{\sin((N+\frac12)t)}{\sin(t/2)}$.\\
Then $s_N(f;x)=\frac{1}{2\pi}\int_{-\pi}^{\pi}f(x-t)D_N(t)\,dt$.\\
If, for some $x$, $\abs{f(x+t)-f(x)}\leq M\abs{t}$ whenever $\abs{t}<\delta$, then $s_N(f;x)\to f(x)$.\\
Hence $f=g$ near $x$ implies $s_N(f;x)-s_N(g;x)\to0$.
\begin{proof}
The displayed form of $D_N$ is the finite geometric sum; substituting $\widehat f(m)$ gives the convolution formula.\\
Also $\frac1{2\pi}\int_{-\pi}^{\pi}D_N=1$. Set $q_x(t)\defeq \frac{f(x-t)-f(x)}{\sin(t/2)}$, for $0<\abs t\leq\pi$,$q_x(0)\defeq0$.\\
Then $s_N(f;x)-f(x)=\frac{1}{2\pi}\int_{-\pi}^{\pi} q_x(t)\sin\!\left(N+\frac12\right)t\,dt$.\\
The hypothesis makes $q_x$ bounded near $0$; away from $0$ it is Riemann integrable.\\
Since $\sin((N+\frac12)t)=\sin Nt\cos(t/2)+\cos Nt\sin(t/2)$, Theorem \ref{thm:orthonormal-fourier-best-approximation-bessel} sends the two Fourier coefficients to $0$.\\
Apply the result to $f-g$ for localization.\\
\hlt{Structural Commentary}: The Dirichlet kernel rewrites Fourier partial sums as a local oscillatory integral.\\
\hlt{Key Insight}: Pointwise convergence under a local Lipschitz condition depends only on nearby values of $f$.
\end{proof}
\end{theorem}

\begin{corollary} \label{cor:trigonometric-polynomial-weierstrass} \hlt{Uniform Approximation by Trigonometric Polynomials}\\
If $f\in C(\R)$ is $2\pi$-periodic, then for every $\varepsilon>0$ there is a trigonometric polynomial $P$ such that $\norm{f-P}_\infty<\varepsilon$.
\begin{proof}
Let $\mathbb{T}\defeq\{z\in\C:\abs z=1\}$ and set $F(e^{ix})\defeq f(x)$.\\
The algebra $\mathcal{T}\defeq\operatorname{span}\{z^n:n\in\Z\}\subset C(\mathbb{T})$ is self-adjoint, separates points, and vanishes nowhere.\\
By Theorem \ref{thm:complex-stone-weierstrass}, $\overline{\mathcal{T}}=C(\mathbb{T})$. Pulling approximants back by $z=e^{ix}$ gives the desired $P$.\\
\hlt{Structural Commentary}: Periodic approximation is Stone--Weierstrass on the unit circle.\\
\hlt{Key Insight}: The functions $e^{inx}$ are the algebraic coordinates on the circle.
\end{proof}
\end{corollary}

\begin{theorem} \label{thm:parseval-fourier-series} \hlt{Parseval's Theorem for Trigonometric Fourier Series}\\
Let $f,g$ be Riemann integrable and $2\pi$-periodic. Let $\langle u,v\rangle_2\defeq\frac1{2\pi}\int_{-\pi}^{\pi}u(x)\overline{v(x)}\,dx$ and $\norm{u}_2^2\defeq\langle u,u\rangle_2$.\\
If $f\sim_{\mathrm{coef}}\sum_{n\in\Z}\widehat{f}(n)e^{inx}$, and $g\sim_{\mathrm{coef}}\sum_{n\in\Z}\widehat{g}(n)e^{inx}$,\\
then $\norm{f-s_N(f)}_2\to0$; $\langle f,g\rangle_2=\sum_{n\in\Z}\widehat{f}(n)\overline{\widehat{g}(n)}$; and $\norm{f}_2^2=\sum_{n\in\Z}\abs{\widehat{f}(n)}^2$.
\begin{proof}
If $f$ is continuous, Corollary \ref{cor:trigonometric-polynomial-weierstrass} gives trigonometric polynomials $P$ with $\norm{f-P}_2$ arbitrarily small.\\
If $P$ has degree $N_0$, then Theorem \ref{thm:orthonormal-fourier-best-approximation-bessel} gives $\norm{f-s_N(f)}_2\leq\norm{f-P}_2$ for $N\geq N_0$.\\
For Riemann integrable $f$, approximate $f$ in $\norm{\cdot}_2$ by a continuous $2\pi$-periodic $h$.\\
Then $\norm{f-s_N(f)}_2 \leq \norm{f-h}_2+\norm{h-s_N(h)}_2+\norm{s_N(h-f)}_2$ and $\norm{s_N(h-f)}_2\leq\norm{h-f}_2$ (Theorem \ref{thm:orthonormal-fourier-best-approximation-bessel}).\\
Hence $s_N(f)\to f$ in $\norm{\cdot}_2$. Finally, $\langle s_N(f),g\rangle_2 =\sum\limits_{\abs{n}\leq N}\widehat{f}(n)\overline{\widehat{g}(n)}$.\\
Let $N\to\infty$ and use Schwarz' inequality; the last identity is $g=f$.\\
\hlt{Structural Commentary}: Density upgrades Bessel to Parseval.\\
\hlt{Key Insight}: Bessel gives one inequality; density of trigonometric polynomials gives equality.
\end{proof}
\end{theorem}

\begin{method} \label{meth:fourier-structure-proof-patterns} \hlt{Fourier Proof Patterns}
\begin{enumerate}[label=(\roman*)]
\setlength{\itemsep}{0pt}
\item \hlt{Characters}: use orthogonality of $e_n$ to compute coefficients and finite-dimensional projections.
\item \hlt{Kernels}: rewrite partial sums or averages as convolutions with Dirichlet or Fejer kernels.
\item \hlt{Approximation}: use Fejer positivity or Stone--Weierstrass to approximate continuous periodic functions by trigonometric polynomials.
\item \hlt{$L^2$ control}: project onto finite orthonormal spans, then pass to limits by density.
\item \hlt{Localization}: express the error as an oscillatory integral and use local regularity near the point.
\end{enumerate}
\hlt{Key Insight}: Fourier theory is finite-dimensional orthogonal projection plus an approximation theorem.
\end{method}
